\chapter{標本化定理}
    \section{基本形}
        \label{標本化定理の基本形}
        \begin{shadebox}
            $d\in\naturalNumbers,\;W_l>0\;(l=1,2,,\dots,d),\;\Omega\coloneq\prod_{l=1}^d [-W_l,W_l]$とする。
            $f:\realNumbers^d\to\complexNumbers$のFourier変換$\FT{f}$が存在してその台が$\Omega$に含まれるとき，次式が成り立つ。
            \[ f(\bm{x}) = \sum_{\bm{n}\in\integers^d}f\left(\pi\frac{n_1}{W_1},\dots,\pi\frac{n_d}{W_d}\right)\prod_{l=1}^d\sinc W_l\left(x_l - \pi\frac{n_l}{W_l}\right) \]
            つまり$f$の各点での評価値を沢山集めて$f$を任意の精度で近似できる。
            \par
            角周波数$W_l$のかわりに周波数$F_l=W_l/(2\pi)$を使うと上式は次式になる。
            \begin{equation}
                f(\bm{x}) = \sum_{\bm{n}\in\integers^d}f\left(\frac{n_1}{2F_1},\dots,\frac{n_d}{2F_d}\right)\prod_{l=1}^d\sinc 2\pi F_l\left(x_l - \frac{n_l}{2F_l}\right)
                \label{equation:標本化定理の基本形の周波数版}
            \end{equation}
        \end{shadebox}
        \begin{proof}
            \quad\par
            $\FT{f}$の台が超直方体$\Omega$に含まれるから$\FT{f}$はFourier級数展開できる。
            第 $\bm{n}$ Fourier 係数を$c(\FT{f},\bm{n})$とすると
            \[ \FT{f}(\bm{\omega}) = \sum_{\bm{n}\in\integers^d} c(\FT{f},\bm{n}) \exp i\sum_{l=1}^d n_l\frac{\omega_l}{W_l}\pi \]
            となる。$c(\FT{f},\bm{n})$は次式で求まる。
            \begin{align*}
                c(\FT{f},\bm{n}) &= \left(\prod_{l=1}^d 2W_l\right)^{-1} \integrate{\Omega}{}{\mathcal{F}(f,\bm{\xi}) \exp (-i)\sum_{l=1}^d n_l\frac{\xi_l}{W_l}\pi}{}{\bm{\xi}} \\
                &= (2\pi)^{d/2} \left(\prod_{l=1}^d 2W_l\right)^{-1} (2\pi)^{-d/2} \integrate{\textcolor{darkpastelgreen}{\realNumbers^d}}{}{\mathcal{F}(f,\bm{\xi}) \exp i\sum_{l=1}^d \left(\frac{-n_l}{W_l}\pi\right)\xi_l}{}{\bm{\xi}} \\
                &= (2\pi)^{d/2} \left(\prod_{l=1}^d 2W_l\right)^{-1} \mathcal{F}^{-1}\left(\FT{f},-\pi\bm{n}\HadamardDiv\bm{W}\right) \\
                &= (2\pi)^{d/2} \left(\prod_{l=1}^d 2W_l\right)^{-1} f\left(-\pi\bm{n}\HadamardDiv\bm{W}\right)
            \end{align*}
            $f$は$\FT{f}$のFourier逆変換で次のようにして求まる。
            \begin{align*}
                f(\bm{x}) &= \mathcal{F}^{-1}(\FT{f})(\bm{x}) = (2\pi)^{-d/2} \integrate{\realNumbers^d}{}{\FT{f}(\bm{\omega})\exp i\bm{\omega}^\top\bm{x}}{}{\bm{\omega}} = (2\pi)^{-d/2} \integrate{\Omega}{}{\FT{f}(\bm{\omega})\exp i\bm{\omega}^\top\bm{x}}{}{\bm{\omega}} \\
                &= (2\pi)^{-d/2} \integrate{\Omega}{}{\sum_{\bm{n}\in\integers^d} c(\FT{f},\bm{n}) \left(\exp i\sum_{l=1}^d n_l\frac{\omega_l}{W_l}\pi\right) \exp i\bm{\omega}^\top\bm{x}}{}{\bm{\omega}} \\
                &= (2\pi)^{-d/2} \sum_{\bm{n}\in\integers^d} \integrate{\Omega}{}{c(\FT{f},\bm{n}) \exp i \bm{\omega}^\top\left(\bm{x} + \pi\bm{n}\HadamardDiv\bm{W}\right) }{}{\bm{\omega}} \\
                &= (2\pi)^{-d/2} \sum_{\bm{n}\in\integers^d} \integrate{\Omega}{}{(2\pi)^{d/2} \left(\prod_{l=1}^d 2W_l\right)^{-1} f\left(-\pi\bm{n}\HadamardDiv\bm{W}\right) \exp i \bm{\omega}^\top\left(\bm{x} + \pi\bm{n}\HadamardDiv\bm{W}\right) }{}{\bm{\omega}} \\
                &= \left(\prod_{l=1}^d 2W_l\right)^{-1} \sum_{\bm{n}\in\integers^d} f\left(-\pi\bm{n}\HadamardDiv\bm{W}\right) \integrate{\Omega}{}{\exp i \bm{\omega}^\top\left(\bm{x} + \pi\bm{n}\HadamardDiv\bm{W}\right) }{}{\bm{\omega}}
            \end{align*}
            ここで
            \begin{align*}
                \integrate{\Omega}{}{\exp i \bm{\omega}^\top\left(\bm{x} + \pi\bm{n}\HadamardDiv\bm{W}\right) }{}{\bm{\omega}} &= \prod_{l=1}^d \integrate{-W_l}{W_l}{\exp i\left(x_l + \pi\frac{n_l}{W_l}\right)\omega_l}{}{\omega_l} \\
                &= \prod_{l=1}^d \frac{1}{i\left(x_l + \pi\frac{n_l}{W_l}\right)}\left[\exp i\left(x_l + \pi\frac{n_l}{W_l}\right)W_l - \exp (-i)\left(x_l + \pi\frac{n_l}{W_l}\right)W_l\right] \\
                &= \prod_{l=1}^d 2W_l \frac{\sin \left(x_l + \pi\frac{n_l}{W_l}\right)W_l}{\left(x_l + \pi\frac{n_l}{W_l}\right)W_l} = \prod_{l=1}^d 2W_l \prod_{l=1}^d \sinc W_l\left(x_l + \pi\frac{n_l}{W_l}\right)
            \end{align*}
            であるから
            \begin{align*}
                f(\bm{x}) &= \left(\prod_{l=1}^d 2W_l\right)^{-1}\sum_{\bm{n}\in\integers^d} f\left(-\pi\bm{n}\HadamardDiv\bm{W}\right) \prod_{l=1}^d 2W_l \sinc \left(x_l + \pi\frac{n_l}{W_l}\right)W_l \\
                &= \sum_{\bm{n}\in\integers^d} f\left(\pi\bm{n}\HadamardDiv\bm{W}\right) \prod_{l=1}^d \sinc W_l\left(x_l - \pi\frac{n_l}{W_l}\right) \\
                &= \sum_{\bm{n}\in\integers^d} f\left(\pi\frac{n_1}{W_1},\dots,\pi\frac{n_d}{W_d}\right) \prod_{l=1}^d \sinc W_l\left(x_l - \pi\frac{n_l}{W_l}\right)
            \end{align*}
        \end{proof}
    \section{再標本化への応用}
        \newcommand*{\fsampWithSubsc}[1]{f_{\text{s},#1}}
        \subsection{動機}
            実用の場面では，あるサンプリング周波数 $\fsampWithSubsc{1}$ で標本化された離散時間信号 $x:\integers\to\complexNumbers$ を別のサンプリング周波数 $\fsampWithSubsc{2}$ で再標本化することがある。
            その際に \cref{equation:標本化定理の基本形の周波数版} を直接適用すると，補間係数すなわち sinc 関数の部分が十分小さくなって計算誤差が許容できる程度に至るまで部分和の計算が必要である。
            $x$ の占有帯域が $[-\fsampWithSubsc{2}/2,\fsampWithSubsc{2}/2]$ 一杯に広がっている場合はそうするより他にない。
            しかし特に通信分野に於いてはその仕様は稀であり，普通はその 80 \% 以下と規定されることが多い\footnote{例えば Analog Devices Inc. が販売しているような DAC, ADC と DSP を組み合わせたデバイスのレート変換処理に於けるフィルタの仕様はそうである。}。
            その場合は，補間係数を sinc 関数より早く 0 に収束する別の係数で置き換えることができる。
            そのような係数は，feedforward フィルタの係数を得る方法として広く用いられている，Remez アルゴリズムで数値的に計算される場合が殆どである\footnote{Raised-Cosine フィルタでもよい。それなら解析的に求められる}。
        \subsection{導出}
            本質に焦点を当てるため 1 次元の場合 ($d=1$) を考えるが，$d>1$ への拡張は容易である。